% ══════════════════════════════════════════════════════════════════════════════
% Chapter 1: Introduction
% ══════════════════════════════════════════════════════════════════════════════

\section{Background}

Person re-identification (Re-ID) is a fundamental computer vision task that involves recognizing and matching individuals across images captured by different, non-overlapping surveillance cameras. Given a query image of a person (or vehicle), the objective is to retrieve all images of the same identity from a large gallery set, regardless of variations in viewpoint, pose, lighting, and background. This task has critical applications in video surveillance, public safety, autonomous driving, and smart retail systems~\cite{zheng2015scalable}.

While supervised Re-ID methods have achieved impressive performance by training on large-scale identity-labelled datasets, the manual annotation of such datasets is prohibitively expensive and time-consuming. This has motivated significant research into \textit{unsupervised} person Re-ID, which aims to learn discriminative feature representations without any identity labels.

\section{Unsupervised Person Re-ID: The Challenge}

Unsupervised Re-ID methods typically follow an iterative two-step framework:
\begin{enumerate}
    \item \textbf{Pseudo label generation:} Features are extracted from a pre-trained backbone, and clustering algorithms (commonly DBSCAN~\cite{ester1996dbscan}) are applied to assign pseudo identity labels to each training sample.
    \item \textbf{Model training:} The model is trained using these noisy pseudo labels via classification and metric learning losses, and the process repeats.
\end{enumerate}

A key challenge in this pipeline is that pseudo labels generated from global features are inherently noisy --- samples from the same person may be assigned to different clusters (under-segmentation), or samples from different persons may be grouped together (over-segmentation). Recent approaches attempt to mitigate this noise by incorporating fine-grained \textit{part-based features} that capture local body region information. However, existing part-based methods often compromise the discriminative semantic structure of the generated patches and overlook the consistency between local and global representations.

\section{The PISL Framework}

To address these limitations, Tao~\textit{et~al.}~\cite{tao2024unsupervised} proposed the \textbf{Person Intrinsic Semantic Learning (PISL)} framework, published in IEEE Transactions on Image Processing (TIP 2024). PISL introduces a novel integration of \textit{diffusion models} into the unsupervised Re-ID pipeline. The core insight is that a diffusion model can generate spatial transformer parameters that extract patches with \textit{intrinsic semantic structure} --- for example, consistently localizing the head, torso, and legs of a person, or the front, side, and rear of a vehicle.

The framework proposes three key innovations:
\begin{itemize}
    \item \textbf{Spatial Diffusion Model (SDM):} Performs a denoising diffusion process from noisy spatial transformer parameters to semantic parameters, enabling the sampling of semantically structured patches.
    \item \textbf{Semantic Controlled Diffusion (SCD) Loss:} Guides the denoising direction of the diffusion model to ensure generated patches capture meaningful semantic regions.
    \item \textbf{Patch Semantic Consistency (PSC) Loss:} Captures consistency between patch and global features, refining pseudo labels through a label refinement mechanism.
\end{itemize}

Additionally, the framework incorporates \textbf{camera-aware contrastive learning}~\cite{wang2021cap} to mitigate the domain gap caused by varying camera conditions across different surveillance viewpoints.

\section{Objective of This Work}

The objective of this thesis is to reproduce and analyse the camera-aware PISL framework on two benchmark datasets:
\begin{itemize}
    \item \textbf{Market-1501}~\cite{zheng2015scalable}: A large-scale person Re-ID benchmark with 1,501 identities across 6 cameras.
    \item \textbf{VeRi-776}~\cite{liu2016deepveri}: A large-scale vehicle Re-ID benchmark with 776 vehicle identities across 20 cameras.
\end{itemize}

Our specific objectives include:
\begin{enumerate}
    \item Faithful reproduction of the PISL framework using the official codebase on available GPU infrastructure.
    \item Quantitative evaluation using standard metrics (mAP and CMC) with comparison against the original paper's results.
    \item Comprehensive qualitative analysis through attention heatmaps, retrieval ranking visualizations, and t-SNE feature embeddings.
    \item Analysis of training dynamics, loss convergence, and factors affecting reproducibility.
\end{enumerate}

\section{Organization of the Report}

The remainder of this report is organized as follows:
\begin{itemize}
    \item \textbf{Chapter 2} reviews the relevant literature on unsupervised person Re-ID, part-based representations, and diffusion models.
    \item \textbf{Chapter 3} presents the detailed methodology of the PISL framework, including the mathematical formulations of all key components.
    \item \textbf{Chapter 4} describes the experimental setup and presents both quantitative and qualitative results.
    \item \textbf{Chapter 5} provides analysis, discussion of the results, conclusions, and directions for future work.
\end{itemize}
